% ----------------------------------------------------------
% Conclusão e Trabalhos Futuros
% Capítulo sem numeração, mas presente no Sumário
% ----------------------------------------------------------

\chapter*[Conclusão e Trabalhos Futuros]{Conclusão e Trabalhos Futuros}\label{cap:conclusao}
\addcontentsline{toc}{chapter}{Conclusão e Trabalhos Futuros}

Este trabalho se propôs a comparar, num mesmo simulador da plataforma F1TENTH e
sob os mesmos critérios de avaliação, três formas de pilotar um veículo de
corrida em escala 1:10: controladores clássicos, políticas obtidas por
aprendizado de máquina e pilotos humanos. O que ele entrega e o que ele não
entrega separam-se por uma fronteira nítida, e é ela que organiza esta
conclusão. Estão medidos, em simulação e sob um protocolo comum, os cinco pilotos
automáticos: o \textit{pure pursuit} do E1, o Disparity Extender e o seguidor de
gap por descontinuidade, do lado clássico, e a rede de imitação do E2 e o resíduo
de aprendizado por reforço do E3, do lado do aprendizado de máquina. Estão
prontos a interface de pilotagem humana e o protocolo da coleta. Faltam as
sessões com os pilotos humanos, e cada ponto desta conclusão que depende delas
está marcado como pendente.

Esta conclusão retoma do Capítulo~\ref{cap:resultados} apenas os números que
sustentam as suas afirmações e não reconstrói o argumento do
Capítulo~\ref{cap:metodologia}. Ela faz duas coisas:
enuncia o que o trabalho sustenta e responde um a um aos quatro objetivos
específicos declarados na seção \textit{Objetivos} da Introdução.

\section*{Conclusões}
\addcontentsline{toc}{section}{Conclusões}

\subsection*{O que o trabalho mediu}

A cascata de três etapas cumpriu, em simulação, o que se propôs a fazer, e a
Tabela~\ref{tab:res_sintese} do Capítulo~\ref{cap:resultados} a resume numa linha
por piloto, ao lado dos demais pilotos do comparativo. O especialista geométrico da etapa E1 fixou o tempo de referência
$T_{pp} = 59{,}000$~s, com 50 de 50 voltas sem colisão e com determinismo
verificado --- o valor reproduziu dígito a dígito uma medição anterior feita em
outra arquitetura, o que é o que autoriza usá-lo como denominador de todos os
critérios seguintes. A rede de imitação da etapa E2 aprovou os quatro critérios
do seu portão e pagou 1,8\% de tempo de volta pela troca que define o trabalho:
deixar de exigir pose global. E o resíduo de aprendizado por reforço da etapa E3
aprovou os seus quatro critérios completando 91 de 92 voltas sob randomização e
percorrendo o circuito em 52,525~s, \textbf{11,0\% mais rápido que o especialista
que originou os seus rótulos}, sem devolver a dependência de localização e com
degradação desprezível sob randomização do coeficiente de atrito --- queda de
taxa de conclusão nula, e piora de tempo de volta de 0,0019 contra um limiar de
0,05.

As três ressalvas que o Capítulo~\ref{cap:resultados} discute nominalmente
acompanham esses números e não são atenuadas em lugar nenhum do texto: o limiar
de taxa de conclusão cai dentro do intervalo de confiança nos três
portões aprovados (Tabela~\ref{tab:val_ic}); a política composta da etapa E3 é
sensível a perturbação numérica irrelevante numa amplitude quatro vezes maior que
a sua margem sobre o limiar, de modo que o número reproduzível ali é a
distribuição e não a realização (Tabela~\ref{tab:val_sensibilidade}); e duas das
quatro pistas do critério de generalização da etapa E2 haviam entrado no conjunto
de coleta, o que dilui aquele critério em vez de inflá-lo
(Seção~\ref{sec:val-heldout}). Duas dessas três estavam declaradas, com o motivo
quantificado, antes de qualquer medição; a terceira foi encontrada conferindo o
artefato contra o texto que o descrevia, e foi declarada em vez de corrigida
retroativamente.

Uma outra ressalva, declarada nos mesmos capítulos fora da lista das três, é de \textbf{protocolo}: o
denominador $T_{pp}$ é medido \textbf{sem} randomização, enquanto as etapas E2 e
E3 são medidas \textbf{com} randomização de atrito, de posição de largada e de
ruído de sensor. A comparação entre 52,525~s e 59,000~s não é, portanto, entre
medições feitas sob o mesmo protocolo. A direção do efeito favorece a leitura
conservadora --- a política que supera a referência é a que corre em condições
mais duras, e o braço sem randomização da comparação pareada da etapa E3 estima o
tamanho disso em cerca de um décimo de segundo ---, mas o padrão deste trabalho é
declarar a assimetria, não confiar na direção dela. O comparativo a mede
diretamente: com randomização e nas mesmas sementes, o especialista fez 59,00~s, e
o resíduo, 52,55~s.

\subsection*{O que o comparativo mediu}

Sob o protocolo comum, os pilotos automáticos se separam em três grupos. Os
pilotos aprendidos e o especialista completam a volta entre 52,55~s e 60,05~s: o
resíduo em 52,55~s, o especialista em 59,00~s, com uma colisão em 45 voltas, e a
rede de imitação em 60,05~s, com taxa de esterço e desvio da linha de corrida
praticamente iguais aos do especialista que ela imita. O Disparity Extender
completa todas as 50 voltas, mas em 145,55~s, com a maior taxa de esterço e o
maior desvio da linha de corrida. E o seguidor de gap não completa volta alguma:
nenhuma das 48 combinações das duas varreduras de parâmetros completou uma volta
de sonda, e a configuração congelada colidiu nas cinco sementes da avaliação. Com
o mesmo sensor em execução, as políticas aprendidas alcançaram o piloto que usa
a pose e a linha de corrida calculada sobre o mapa, o E1; os controladores
reativos, que não usam mapa, não. A ressalva central dessa leitura está na
Seção~\ref{sec:res-disc-cmp}: as políticas aprendidas não usam o mapa em
execução, mas a informação dele entrou no treino, pelas demonstrações do
especialista.

\subsection*{A contribuição metodológica}

O resultado que este trabalho considera mais transferível não é o tempo de volta.
É o aparato descrito no Capítulo~\ref{cap:metodologia}: um contrato numérico de
118 asserções com identificador estável, cada uma registrando a falha medida que a
originou; cinco camadas de verificação que vão da cobertura do contrato até a
proveniência do repositório; critérios de
aceite congelados no nascimento de cada execução, com a proveniência tratada como
parte do veredito e não como metadado; e uma tabela de ações de correção
pré-declaradas por critério, que é o que separa um experimento de uma busca
livre.

Nada disso é caro em si, e é essa a razão de o capítulo ser escrito como
argumento e não como manual. O que custa é a disciplina de aplicar cada regra
antes de a falha aparecer --- e a evidência oferecida é que, neste projeto, cada
uma delas foi escrita \textit{depois} de uma falha medida, cujo custo variou
entre uma tentativa de uma hora e oito execuções reprovadas em silêncio. Quatro
episódios sustentam o argumento inteiro: um detector de colisão que deixava o
veículo atravessar a parede e seguir dirigindo fora do circuito, um vazamento de
estado entre episódios que se disfarçava de má generalização, duas cópias da mesma
composição de política divergindo por um tipo numérico, e três valores distintos
para a mesma medição convivendo no repositório. Nos quatro, reexecutar o comando
devolvia o mesmo número; nos quatro, o número estava errado; e nos quatro a
distância até o valor correto era da mesma ordem das diferenças que o trabalho
pretendia demonstrar.

\subsection*{Os quatro objetivos, um a um}

A Introdução declarou quatro objetivos específicos. Dizer com precisão quais foram
atingidos, quais o foram em parte e quais não o foram é mais forte do que
dissolvê-los numa afirmação geral de êxito --- e é a única conduta compatível
com o Capítulo~\ref{cap:metodologia}, que argumenta o capítulo inteiro que
declarar o que não se cobre é parte do resultado. A
Tabela~\ref{tab:conc_objetivos} resume, e os parágrafos seguintes qualificam cada
linha.

\begin{table}[h]
    \centering
    \small
    \begin{tabularx}{\textwidth}{|c|>{\raggedright\arraybackslash}X|>{\raggedright\arraybackslash}p{3.3cm}|}
        \hline
        \textbf{\#} & \textbf{Objetivo específico declarado} & \textbf{Situação} \\
        \hline
        1 & Implementar e avaliar em simulação três controladores clássicos: o
        especialista \textit{pure pursuit}, que usa a linha de corrida calculada
        sobre o mapa, e dois reativos que usam apenas o LiDAR, o Disparity
        Extender e o seguidor de gap &
        atingido; o seguidor de gap não completa volta \\
        \hline
        2 & Treinar uma política de imitação sem pose global e refiná-la por
        aprendizado por reforço residual, sob critérios fixados antes da
        medição & atingido \\
        \hline
        3 & Desenvolver uma interface de pilotagem humana no mesmo simulador e
        coletar voltas de pilotos humanos sob protocolo fixado antes da coleta
        & em parte: interface e protocolo prontos, coleta pendente \\
        \hline
        4 & Comparar os três tipos de piloto quanto a conclusão de volta, tempo de
        volta, colisões, suavidade do esterço, desvio da raceline e perfil de
        velocidade, com o humano em duas condições de informação & em parte:
        pilotos automáticos comparados, humanos pendentes \\
        \hline
    \end{tabularx}
    \caption{Os quatro objetivos específicos declarados na Introdução diante do
    que o trabalho de fato mediu. A situação dos objetivos 3 e 4 é a desta
    versão, anterior às sessões humanas}
    \label{tab:conc_objetivos}
\end{table}

\textbf{Objetivo 1.} Os três controladores clássicos foram implementados e
avaliados sob o protocolo comum. O \textit{pure pursuit} tem dois números, que
convivem e não se substituem: o oficial, sem randomização, de 59,000~s e 50 de
50 voltas, que é o denominador $T_{pp}$ dos critérios de aceite; e o do
comparativo, com randomização, de 59,00~s e 44 de 45 voltas. O seguidor de gap foi
portado do nó ROS2 para o ambiente de pesquisa, passou pela varredura da grade
oficial e por uma grade ampliada criada depois dela, e foi avaliado na
configuração congelada: 0 de 5 voltas, com cinco colisões. O objetivo pedia
implementar e avaliar, e isso foi feito; o resultado da avaliação é negativo, e a
Seção~\ref{sec:res-disc-cmp} discute as causas prováveis. O Disparity Extender,
portado pelo mesmo caminho, completa 50 de 50 voltas em 145,55~s.

\textbf{Objetivo 2.} A rede de imitação sem pose global (E2) e o resíduo de
aprendizado por reforço (E3) foram treinados e avaliados sob portões cujos
limiares foram congelados antes da medição, e ambos aprovaram os seus critérios,
com as ressalvas declaradas acima (E2: 60,050~s e 50 de 50 voltas; E3: 52,525~s e
91 de 92 voltas, ambos com randomização). Reavaliados sob o protocolo comum, sem
retreino e já sob o avaliador corrigido, reproduziram as medianas a menos de
0,03~s (60,05~s e 52,55~s, ambos com 50 de 50 voltas).
A percepção consumida pela política é exclusivamente LiDAR bidimensional e
proprioceptiva, o que está declarado como limite na Seção~\ref{sec:arte-lacuna}.

\textbf{Objetivo 3.} A interface de pilotagem humana está pronta: \textit{joystick}
com os mesmos limites e o mesmo atraso das políticas, passo segurado em tempo
real com o atraso registrado, duas vistas e gravação no mesmo formato dos pilotos
automáticos (Seção~\ref{sec:prop-humano}). O protocolo da coleta está escrito,
com \textit{hash} gravado em cada registro (Seção~\ref{sec:val-humano}). A coleta
não foi feita: \pendentehumano{pilotos coletados (autores e voluntários), sessões
completas e excluídas, e a situação do Comitê de Ética}. A participação de
voluntários depende da resposta sobre a necessidade de submissão ao Comitê de
Ética em Pesquisa; sem ela, o estudo se limita aos autores.

\textbf{Objetivo 4.} A comparação entre os pilotos automáticos está feita, nas
seis métricas do objetivo, e a Tabela~\ref{tab:comparativo} já reserva as linhas
humanas. Falta a parte humana: \pendentehumano{faixa do grupo humano por vista e
posição de cada política em relação a ela}. Das limitações declaradas antes da
coleta (Seção~\ref{sec:val-limitacoes-comparativo}), cinco pesam diretamente
sobre essa parte: a informação desigual na vista da pista inteira, a ordem fixa
das vistas, a amostra pequena, a separação entre autores e voluntários e a pista
única.

\subsection*{O que fica dito com precisão}

Dos quatro objetivos, dois estão atingidos e dois em parte, e nenhum ficou sem
execução: os dois incompletos dependem da mesma coleta. O que o trabalho afirma,
nesta versão, é o que está medido: uma política que enxerga apenas LiDAR e
proprioceptiva completa 91 de 92 voltas sob randomização no seu portão e percorre
Spielberg 11,0\% mais rápido que o especialista geométrico que a ensinou; sob o
protocolo comum, ela faz 52,55~s contra 59,00~s do especialista, e os dois
controladores reativos que usam o mesmo sensor ficam muito atrás ou não
completam a volta. A afirmação sobre os humanos fica para depois da coleta:
\pendentehumano{afirmação central sobre humano e máquina, por vista}.

\section*{Trabalhos Futuros}
\addcontentsline{toc}{section}{Trabalhos Futuros}

Parte do que resta é continuação do próprio comparativo; parte são direções que a
literatura cobre e este trabalho deixa de fora.

\subsection*{As lacunas do próprio resultado em simulação}

A Seção~\ref{sec:res-discussao} lista sem atenuação o que os números não
sustentam. Cinco itens dali são acionáveis dentro do próprio simulador, e valem
como trabalho futuro imediato:

\begin{itemize}
    \item \textbf{Medir a variabilidade entre sementes de treino.} Cada
          configuração da etapa E3 corresponde a \textit{um} treino, e as
          sementes dos portões são de avaliação. Sem isso, não se sabe quanto da
          diferença entre a configuração aprovada e a anterior cabe dentro da
          dispersão de dois treinos idênticos. É a lacuna mais direta a fechar.
    \item \textbf{Avaliar o resíduo fora do circuito em que treinou.} Todos os
          critérios de aceite das três etapas são medidos em Spielberg, e a etapa
          E3 nunca foi pontuada em nenhuma outra pista: o valor de generalização
          do resíduo simplesmente não existe.
    \item \textbf{Reexecutar a etapa E2 sob a correção da sobreposição} entre o
          conjunto de coleta e o de generalização. A correção já está congelada
          no contrato (\texttt{SPEC-GATE-09}): nenhuma ação pré-declarada pode
          coletar numa pista de generalização. A
          Seção~\ref{sec:val-heldout} explica por que ela não podia ser
          aplicada retroativamente, e a execução reportada segue parcialmente
          contaminada até ser substituída.
    \item \textbf{Refazer o estudo de sensibilidade numérica nesta máquina.} As
          execuções que o sustentam são da máquina anterior do projeto e não
          sobreviveram à migração, de modo que a confrontação entre a margem da
          etapa E3 e a amplitude sob perturbação combina um número auditável com
          outro que já não é reconferível.
    \item \textbf{Tornar a evidência durável e completar o aparato.} Os
          relatórios das execuções oficiais não são versionados e o relatório
          de verificação é sobrescrito a cada execução; promovê-los para um
          caminho versionado, como já foi feito com os relatórios das cinco
          execuções do comparativo, torna toda medição reconferível. Os pontos
          de verificação das três etapas já foram medidos sob o avaliador
          corrigido, em que a colisão tem precedência sobre a chegada, no
          comparativo (Seção~\ref{sec:res-cmp-reavaliacao}); falta refazer sob
          ele os portões oficiais, com os seus critérios, e corrigir as poses de
          largada que nascem em colisão em Monza e Oschersleben, aceitando a
          perda de comparabilidade com as execuções anteriores.
\end{itemize}

\subsection*{Extensões do comparativo}

Do próprio comparativo, quatro extensões são diretas:

\begin{itemize}
    \item \textbf{Validação em hardware.} Todos os pilotos deste trabalho correm
          em simulação. Repetir o comparativo num veículo físico é a extensão de
          maior custo, e este trabalho não a aborda.
    \item \textbf{Imitação de pilotos humanos.} As voltas humanas gravadas passo a
          passo, no mesmo formato das demonstrações do especialista, podem servir
          de demonstração para a rede de imitação, no lugar do especialista
          geométrico.
    \item \textbf{Outras pistas.} O comparativo usa só Spielberg. Em pistas não
          vistas, a independência de mapa dos controladores reativos passa a
          contar, e a generalização das políticas aprendidas e dos humanos pode ser
          medida.
    \item \textbf{Contrabalanço das vistas.} A ordem fixa das vistas mistura o
          efeito da informação com o da prática. Alternar a ordem entre pilotos,
          com amostra maior, separa os dois.
\end{itemize}

Além dessas, permanecem em aberto as direções que a Seção~\ref{sec:arte-lacuna}
identifica como ausentes deste trabalho e presentes na literatura: corrida com
dois veículos e ultrapassagem, adaptação de modelo em linha, garantias formais
sobre a política e percepção além do LiDAR bidimensional.

\subsection*{Consideração final}

O trabalho termina com um comparativo em simulação cuja parte automática está
medida e cuja parte humana depende de coleta ainda por fazer. Sob o mesmo
protocolo, uma política que só vê o LiDAR, obtida por imitação de um especialista
geométrico e refinada por um resíduo de aprendizado por reforço, percorre
Spielberg mais rápido que esse especialista; o controlador reativo que usa o mesmo
sensor sem treino completa a volta em 145,55~s, e o outro não a completa. Onde os pilotos humanos se colocam diante disso,
\pendentehumano{uma frase com o resultado humano}. Dizer isso nessa ordem, com o
que está medido separado do que falta medir, é a mesma prática que o
Capítulo~\ref{cap:metodologia} defende como contribuição: declarar o que se
cobre, declarar o que não se cobre, e quantificar a distância entre os dois sempre
que houver medição que permita.
